\ejercicio{Investigaci\'on: Clase \texttt{Scanner} en Java}

Investigue la clase \texttt{Scanner} en el lenguaje Java, e int\'egrela en un ejercicio que lo aplique.

\subsection*{Documentaci\'on oficial}

Seg\'un la documentaci\'on oficial de Oracle (Java SE 17):

\begin{quote}
\textit{``A simple text scanner which can parse primitive types and strings using
regular expressions. A Scanner breaks its input into tokens using a delimiter
pattern, which by default matches whitespace. The resulting tokens may then be
converted into values of different types using the various \texttt{next} methods.''}
\end{quote}

\textbf{Fuente:} \href{https://docs.oracle.com/en/java/javase/17/docs/api/java.base/java/util/Scanner.html}{\texttt{docs.oracle.com --- java.util.Scanner}}

\subsection*{Caracter\'isticas principales}

\begin{itemize}[itemsep=0.3em]
    \item Pertenece al paquete \texttt{java.util} e implementa las interfaces
          \texttt{Iterator<String>} y \texttt{Closeable}.
    \item Se instancia apuntando a una fuente de entrada:
          \texttt{Scanner sc = new Scanner(System.in);}
    \item Principales m\'etodos de lectura:
    \begin{itemize}[itemsep=0.2em]
        \item \texttt{nextLine()} --- lee una l\'inea completa como \texttt{String}
        \item \texttt{next()} --- lee una sola palabra (hasta el espacio)
        \item \texttt{nextInt()} --- lee un n\'umero entero
        \item \texttt{nextDouble()} --- lee un n\'umero decimal
        \item \texttt{hasNextInt()} --- verifica si el pr\'oximo token es un entero v\'alido
    \end{itemize}
    \item Siempre debe cerrarse al terminar: \texttt{sc.close();}
\end{itemize}

\subsection*{Advertencia: problema del buffer}

Al mezclar \texttt{nextInt()} o \texttt{nextDouble()} con \texttt{nextLine()}, el salto de
l\'inea (\texttt{\textbackslash n}) queda en el buffer y la siguiente lectura de texto
resulta vac\'ia. La soluci\'on es agregar un \texttt{scanner.nextLine()} de limpieza
inmediatamente despu\'es de leer un n\'umero.

\newpage
\subsection*{Aplicaci\'on pr\'actica: Adivina el n\'umero}

Se desarrolla un juego interactivo de consola donde el usuario dispone de 7 intentos
para adivinar un n\'umero secreto. El ejercicio aplica \texttt{nextInt()} para leer
los intentos del usuario y \texttt{hasNextInt()} para validar que la entrada sea un
entero antes de procesarla, evitando una excepci\'on \texttt{InputMismatchException}.

\codigojava{ejercicios/ej10/Ejercicio10.java}

\vspace{1em}
\textbf{Salida del programa:}
\vspace{0.5em}
\begin{figure}[H]
    \centering
    \includegraphics[width=1\textwidth]{ejercicios/ej10/captura.png}
\end{figure}

\volverindice
\newpage